\documentclass[aspectratio=169, 10pt]{beamer}

\usetheme{Madrid}
\usefonttheme{professionalfonts}

\usepackage[utf8]{inputenc}
\usepackage[indonesian]{babel}
\usepackage{amsmath,amsfonts,amssymb}
\usepackage{graphicx}
\usepackage{booktabs}
\usepackage{listings}
\usepackage{xcolor}

\definecolor{mainblue}{RGB}{24, 43, 73}
\definecolor{accentblue}{RGB}{41, 128, 185}
\definecolor{accentgold}{RGB}{230, 126, 34}
\definecolor{darkgray}{RGB}{44, 62, 80}
\definecolor{lightbg}{RGB}{248, 250, 252}
\definecolor{bordergray}{RGB}{200, 210, 220}
\definecolor{codebg}{RGB}{240, 243, 246}
\definecolor{codegreen}{RGB}{39, 174, 96}

\setbeamercolor{structure}{fg=mainblue}
\setbeamercolor{palette primary}{bg=mainblue, fg=white}
\setbeamercolor{palette secondary}{bg=accentblue, fg=white}
\setbeamercolor{palette tertiary}{bg=mainblue!80, fg=white}
\setbeamercolor{title}{fg=white, bg=mainblue}
\setbeamercolor{frametitle}{fg=white, bg=mainblue}
\setbeamercolor{block title}{bg=mainblue, fg=white}
\setbeamercolor{block body}{bg=lightbg, fg=darkgray}
\setbeamercolor{block title alerted}{bg=accentgold, fg=white}
\setbeamercolor{block body alerted}{bg=accentgold!10, fg=darkgray}

\lstset{
    language=Python,
    backgroundcolor=\color{codebg},
    basicstyle=\ttfamily\scriptsize\color{darkgray},
    keywordstyle=\bfseries\color{accentblue},
    stringstyle=\color{accentgold},
    commentstyle=\itshape\color{codegreen},
    showstringspaces=false,
    breaklines=true,
    frame=single,
    rulecolor=\color{bordergray},
    numbers=none
}

\title[\textbf{Deret Waktu (Minggu 4)}]{\textbf{Pemodelan Stokastik Linier: AR, MA, dan ARMA}}
\subtitle{Representasi Operator Polinomial, Analisis Kestabilan Bidang Z, dan Kriteria AIC/BIC}
\author[\textbf{Ir. Novalio Daratha, Ph.D.}]{\textbf{Ir. Novalio Daratha, S.T., M.Sc., Ph.D.}}
\institute[FT UNIB]{Program Studi Teknik Elektro -- Fakultas Teknik\\Universitas Bengkulu}
\date{Minggu 4 / Semester Ganjil 2026}

\begin{document}

\begin{frame}[plain]
    \titlepage
\end{frame}

\begin{frame}{Agenda Perkuliahan Minggu ke-4}
    \begin{columns}[T]
        \begin{column}{0.48\textwidth}
            \begin{block}{Capaian Pembelajaran (Sub-CPMK 4)}
                Mahasiswa mampu menyusun formulasi matematis model AR($p$), MA($q$), dan ARMA($p,q$), menganalisis kestabilan dan invertibilitas pada bidang Z, serta memilih ordo optimal menggunakan kriteria informasi AIC/BIC.
            \end{block}
        \end{column}
        \begin{column}{0.48\textwidth}
            \begin{enumerate}
                \item \textbf{Model AutoRegressive (AR($p$))}
                \item \textbf{Model Moving Average (MA($q$))}
                \item \textbf{Model Gabungan ARMA($p,q$)}
                \item \textbf{Kestabilan \& Bidang Z (Karakteristik Fisik)}
                \item \textbf{Kriteria Informasi: AIC, AICc, dan BIC}
                \item \textbf{Diagnostik Residu \& Implementasi Python}
            \end{enumerate}
        \end{column}
    \end{columns}
\end{frame}

\begin{frame}{Model AutoRegressive (AR($p$))}
    Nilai sekarang $y_t$ diregresikan terhadap nilai masa lalu $y_{t-1}, \dots, y_{t-p}$:
    \begin{block}{Persamaan AR($p$)}
        \begin{equation}
            y_t = c + \phi_1 y_{t-1} + \phi_2 y_{t-2} + \dots + \phi_p y_{t-p} + \epsilon_t \quad \text{di mana } \epsilon_t \sim \mathcal{WN}(0, \sigma^2)
        \end{equation}
    \end{block}
    \begin{columns}[T]
        \begin{column}{0.48\textwidth}
            \textbf{Notasi Operator Lag ($B$):}
            \begin{equation}
                (1 - \phi_1 B - \phi_2 B^2 - \dots - \phi_p B^p)y_t = c + \epsilon_t
            \end{equation}
            $\Phi(B)y_t = c + \epsilon_t$
        \end{column}
        \begin{column}{0.48\textwidth}
            \textbf{Pola Correlogram AR($p$):}
            \begin{itemize}
                \item \textbf{ACF}: Meluruh secara eksponensial (\textit{tail off}).
                \item \textbf{PACF}: Terpotong tajam (\textit{cut off}) tepat setelah $lag = p$.
            \end{itemize}
        \end{column}
    \end{columns}
\end{frame}

\begin{frame}{Estimasi Parameter AR($p$) Murni: Persamaan Yule-Walker}
    Untuk model AR($p$) murni, parameter $\phi_1, \phi_2, \dots, \phi_p$ dapat diestimasi secara deterministik menggunakan autokorelasi sampel ($\rho_k$) melalui sistem persamaan \textbf{Yule-Walker}:
    \begin{block}{Bentuk Matriks Yule-Walker (Sistem Toeplitz)}
        \begin{equation}
            \begin{bmatrix}
                1 & \rho_1 & \rho_2 & \dots & \rho_{p-1} \\
                \rho_1 & 1 & \rho_1 & \dots & \rho_{p-2} \\
                \rho_2 & \rho_1 & 1 & \dots & \rho_{p-3} \\
                \vdots & \vdots & \vdots & \ddots & \vdots \\
                \rho_{p-1} & \rho_{p-2} & \rho_{p-3} & \dots & 1
            \end{bmatrix}
            \begin{bmatrix}
                \phi_1 \\ \phi_2 \\ \phi_3 \\ \vdots \\ \phi_p
            \end{bmatrix}
            =
            \begin{bmatrix}
                \rho_1 \\ \rho_2 \\ \rho_3 \\ \vdots \\ \rho_p
            \end{bmatrix}
            \quad \Longrightarrow \quad \mathbf{R} \boldsymbol{\phi} = \mathbf{r}
        \end{equation}
    \end{block}
    \vspace{0.1cm}
    \begin{itemize}
        \item Solusi analitik didapatkan via Inversi Matriks: $\boldsymbol{\phi} = \mathbf{R}^{-1} \mathbf{r}$
        \item Estimasi variansi \textit{white noise}: $\sigma^2 = \gamma_0 \left( 1 - \sum_{i=1}^p \phi_i \rho_i \right)$
    \end{itemize}
\end{frame}

\begin{frame}{Model Moving Average (MA($q$))}
    Nilai sekarang $y_t$ dimodelkan sebagai kombinasi linier dari kesalahan acak masa lalu:
    \begin{block}{Persamaan MA($q$)}
        \begin{equation}
            y_t = \mu + \epsilon_t + \theta_1 \epsilon_{t-1} + \theta_2 \epsilon_{t-2} + \dots + \theta_q \epsilon_{t-q}
        \end{equation}
        Dalam operator lag: $y_t = \mu + \Theta(B)\epsilon_t$ dengan $\Theta(B) = 1 + \theta_1 B + \dots + \theta_q B^q$.
    \end{block}
    \begin{columns}[T]
        \begin{column}{0.48\textwidth}
            \textbf{Pola Correlogram MA($q$):}
            \begin{itemize}
                \item \textbf{ACF}: Terpotong tajam (\textit{cut off}) setelah $lag = q$.
                \item \textbf{PACF}: Meluruh lambat (\textit{tail off}).
            \end{itemize}
        \end{column}
        \begin{column}{0.48\textwidth}
            \textbf{Invertibilitas MA:}
            Memastikan derau acak masa lalu dapat diestimasi dari observasi sinyal nyata yang terukur.
        \end{column}
    \end{columns}
\end{frame}

\begin{frame}{Kestabilan Sistem Dinamis pada Bidang Z}
    Sebagai insinyur elektro, model ARMA identik dengan filter IIR (Infinite Impulse Response):
    \begin{block}{Syarat Stasioneritas \& Kestabilan}
        Akar-akar polinomial karakteristik $\Phi(B) = 0$ harus berada \textbf{di luar lingkaran satuan} ($|B_i| > 1$) pada bidang $B$, atau akar-akar pada bidang Z ($z = B^{-1}$) harus berada \textbf{di dalam lingkaran satuan} ($|z_i| < 1$):
        \begin{equation}
            1 - \phi_1 z^{-1} - \phi_2 z^{-2} - \dots - \phi_p z^{-p} = 0 \quad \Longrightarrow \quad |z_i| < 1
        \end{equation}
    \end{block}
    \begin{alertblock}{Contoh Kasus AR(1): $y_t = \phi_1 y_{t-1} + \epsilon_t$}
        \begin{itemize}
            \item Jika $|\phi_1| < 1$: Sistem stabil dan stasioner.
            \item Jika $|\phi_1| = 1$: Sistem \textit{Random Walk} (Non-stasioner/Marginally Stable).
            \item Jika $|\phi_1| > 1$: Sistem tak stabil (*Explosive System*).
        \end{itemize}
    \end{alertblock}
\end{frame}

\begin{frame}{Kriteria Informasi untuk Seleksi Ordo (AIC \& BIC)}
    Untuk memilih kombinasi ordo $(p, q)$ terbaik tanpa terjadi \textit{overfitting}:
    \begin{columns}[T]
        \begin{column}{0.48\textwidth}
            \begin{block}{\textbf{Akaike Information Criterion (AIC)}}
                \begin{equation}
                    \text{AIC} = -2\ln(\hat{L}) + 2k
                \end{equation}
                $k = p + q + 1$ (jumlah parameter terestimasi), $\hat{L}$ adalah nilai maksimum fungsi *likelihood*.
            \end{block}
        \end{column}
        \begin{column}{0.48\textwidth}
            \begin{block}{\textbf{Bayesian Information Criterion (BIC)}}
                \begin{equation}
                    \text{BIC} = -2\ln(\hat{L}) + k\ln(N)
                \end{equation}
                Memberikan penalti lebih ketat untuk model yang terlalu kompleks pada dataset besar ($N$).
            \end{block}
        \end{column}
    \end{columns}
    \vspace{0.2cm}
    \begin{center}
        \textbf{Kaidah Keputusan}: Pilih model dengan nilai AIC / BIC terkecil!
    \end{center}
\end{frame}

\begin{frame}[fragile]{Implementasi Python: Estimasi Model ARMA}
\begin{lstlisting}
from statsmodels.tsa.arima.model import ARIMA
import pandas as pd

# Load data stasioner (misal sinyal getaran motor)
df = pd.read_csv('vibration_stationary.csv')

# Grid Search ordo (p, q) berdasarkan nilai AIC minimum
best_aic, best_order = float("inf"), None
for p in range(1, 4):
    for q in range(1, 4):
        try:
            model = ARIMA(df['signal'], order=(p, 0, q)).fit()
            if model.aic < best_aic:
                best_aic = model.aic
                best_order = (p, q)
        except:
            continue

print(f"Model Terbaik: ARMA{best_order} dengan AIC: {best_aic:.2f}")

# Fitting model optimal & tampilkan akar karakteristik
final_model = ARIMA(df['signal'], order=(best_order[0], 0, best_order[1])).fit()
print(final_model.summary())
\end{lstlisting}
\end{frame}

\begin{frame}{Rangkuman \& Persiapan Minggu ke-5}
    \begin{block}{Poin Kunci Minggu 4}
        \begin{itemize}
            \item Model AR mengaitkan sinyal masa kini dengan histori masa lalu; model MA mengaitkan dengan histori galat derau.
            \item Kestabilan ARMA identik dengan kestabilan filter digital IIR (kutub di dalam lingkaran satuan bidang Z).
            \item Nilai AIC dan BIC menjadi kompas penentu ordo parsimonius (*Occam's razor*).
        \end{itemize}
    \end{block}
    \vspace{0.2cm}
    \begin{alertblock}{Materi Minggu Depan (Minggu ke-5)}
        \textbf{Model ARIMA \& SARIMA Box-Jenkins}: Integrasi diferensiasi ($d$) dan komponen musiman $(P,D,Q)_s$ untuk data beban listrik nyata berkala mingguan.
    \end{alertblock}
\end{frame}

\begin{frame}[plain]
    \begin{center}
        \vspace{1.5cm}
        {\Huge \color{mainblue} \textbf{Terima Kasih}}\\[0.4cm]
        {\large Pertanyaan \& Diskusi Terkait Pemodelan ARMA}\\[1.2cm]
        \textbf{Laboratorium Sistem Tenaga \& Komputasi Cerdas}\\
        Jurusan Teknik Elektro -- Universitas Bengkulu
    \end{center}
\end{frame}

\end{document}
